Emergency response operations rarely unfold under stable or fully observable conditions. Dispatchers must make rapid decisions under time pressure, incomplete situational awareness, and dynamically evolving constraints while responders are already en route~\cite{alvear2013}. In such environments, routing, responder allocation, infrastructure readiness, and environmental context cannot be treated as independent components; instead, they must be jointly reasoned in real time to produce actionable and reliable decisions \cite{zografos2002}. Despite advances in intelligent transportation and navigation systems, most existing routing approaches remain fundamentally optimized for nominal travel time, overlooking operational factors such as access limitations, infrastructure availability, and environmental degradation~\cite{jesus2024}. As a result, routes that are computationally optimal may be operationally infeasible or unsafe in real-world emergency scenarios, particularly under dynamic and uncertain conditions.

The increasing availability of Internet of Things (IoT) sensing, distributed infrastructure data, and urban analytics has expanded the potential for data-driven emergency decision support in smart-city environments \cite{asensio2015}. However, integrating these heterogeneous data sources into a unified and operationally meaningful decision-support workflow remains a significant challenge~\cite{park2018}. In practice, emergency vehicle routing is highly sensitive to real-world constraints, including narrow access corridors, gated segments, one-way restrictions, and degraded visibility due to smoke or environmental hazards~\cite{hong2023}. These factors fundamentally alter route feasibility and safety, yet they are often underrepresented or entirely absent in conventional routing formulations \cite{shaaban2019}. As a result, shortest-path solutions may be computationally optimal but operationally impractical~\cite{zhen2024}.

University campuses provide a compelling microcosm of these challenges. Although smaller than urban road networks, campuses exhibit dense, heterogeneous layouts composed of internal roads, service alleys, pedestrian pathways, and restricted-access zones~\cite{sun2024}. In fire response scenarios, the geometrically shortest route is not necessarily the most actionable one. Dispatchers must consider whether routes are traversable by fire vehicles, whether hydrants along the path are operational, and whether environmental conditions compromise maneuverability. Recent work on urban accessibility further emphasizes that operational reachability should be modeled explicitly rather than approximated solely through geometric distance \cite{baik2025}. These observations motivate a shift from purely geometric routing toward operationally decision-support systems.

Despite substantial progress in emergency decision-support systems, three key limitations remain~\cite{sarkar2022}. First, routing is typically treated as an isolated optimization problem, without explicit integration of infrastructure readiness, environmental sensing, and responder state. Second, many systems lack explainability, limiting operator trust and interpretability in high-stakes decision-making scenarios. Third, post-incident accountability and governance are rarely incorporated into system design, resulting in fragmented workflows between dispatch, monitoring, and audit processes. Addressing these limitations requires rethinking emergency routing as part of a broader service-level architecture rather than a standalone algorithmic task. This shift motivates the need for  reproducible systems that can support both real-time decision-making and post-incident analysis within a unified framework.

We present \textit{FireRoute}, a campus-scale emergency routing and dispatch service developed and evaluated on a reproducible pilot testbed around Kasetsart University (Bangkhen). Rather than focusing solely on route optimization, FireRoute frames emergency response as an integrated decision-support workflow that unifies incident intake, explainable risk assessment, risk-weighted routing, hydrant-aware dispatch, responder assignment, and evidence-ready governance. From a big-data service perspective, the system emphasizes the integration of heterogeneous operational data including road constraints, hydrant status, environmental sensing, responder availability, incident descriptors, and audit records into a coherent and inspectable pipeline. The contributions of this paper are as follows:

\begin{enumerate}
    \item We propose a service-oriented formulation of emergency response that integrates routing, environmental sensing, and governance into a decision-support workflow.

    \item We introduce a risk-aware routing and dispatch framework that jointly considers access constraints, infrastructure availability, and environmental conditions, including a hydrant-aware dispatch mechanism that optimizes responder--hydrant--incident chains.

    \item We develop an explainable incident-priority model that exposes interpretable factors influencing dispatch decisions, enabling transparent and trustworthy operation in high-stakes scenarios.

    \item We provide a reproducible system package that externalizes asset, scenario, and governance data, enabling transparent evaluation and reuse for future research.
\end{enumerate}

The remainder of this paper is organized as follows. Section~\ref{sec:background} reviews related work and situates FireRoute within emergency decision-support systems, smart-city sensing, and routing under uncertainty. Section~\ref{sec:methodology} presents FireRoute, including its architecture, data model, and core decision-support mechanisms. Section~\ref{sec:evaluation} describes the evaluation setup, named scenarios, and scenario-based results. Finally, Section~\ref{sec:conclusion} concludes the paper by summarizing the main findings, limitations, and future directions. An unnumbered reproducibility section then summarizes artifact availability and points readers to the detailed execution instructions in the accompanying package.